%----------------------------------------------------------------------------------------
%	Chapter 4: Implementation, Results and Discussions
%----------------------------------------------------------------------------------------
\chapter{Implementation, Results and Discussions}
\label{ch:results}

\section{Introduction}

This chapter presents the implementation details of the HyDART-RQ framework, the experimental setup, the evaluation metrics, and the results across all three datasets. We compare HyDART-RQ against multiple baseline methods and conduct ablation studies to assess the contribution of individual design choices.

\section{Implementation}

The HyDART-RQ framework is implemented entirely in Python 3.12. The key libraries used are:
\begin{itemize}
  \item \textbf{NumPy 1.x}: Vectorised operations, tensor construction (\texttt{numpy.add.at}), partition-based top-$k$ selection.
  \item \textbf{SciPy}: Sparse matrix construction (\texttt{scipy.sparse.csr\_matrix}), truncated SVD (\texttt{scipy.sparse.linalg.svds}).
  \item \textbf{pandas}: Data loading, timestamp parsing, categorical encoding.
  \item \textbf{scikit-learn}: Data partitioning utilities.
  \item \textbf{huggingface\_hub}: Direct access to Amazon Reviews 2023 dataset files.
\end{itemize}

The project is organized as:
\begin{itemize}
  \item \texttt{data/common\_data.py}: Dataset loading and preprocessing functions.
  \item \texttt{baselines/durable/ssa.py}: SSA implementation.
  \item \texttt{baselines/durable/pra.py}: PRA implementation.
  \item \texttt{hybrid/candidate\_filter.py}: DQR and adaptive filter.
  \item \texttt{hybrid/durable\_verifier.py}: Verification on candidate set.
  \item \texttt{utils/rank\_table.py}: Rank table construction.
  \item \texttt{experiments/}: Per-dataset experiment scripts.
  \item \texttt{outputs/}: Experiment results in CSV and TXT format.
\end{itemize}

\subsection{Correctness Validation}

Before benchmarking, we ran eight unit tests comparing HyDART-RQ output against brute-force ground truth on synthetic data with known answers:
\begin{enumerate}
  \item All windows active for all users (durability = 1.0 for all).
  \item No windows active (empty result expected).
  \item Exactly $\tau$-fraction of windows active (boundary case).
  \item One user durable, all others not.
  \item Query item at rank 1 in all windows.
  \item Query item at rank $k+1$ in all windows (no durable users).
  \item Ties in preference scores.
  \item Partial query interval $[t_b, t_e) \subsetneq [0, L)$.
\end{enumerate}
All eight tests pass after the score-convention correction (see Section~\ref{sec:hydart}).

\section{Experimental Setup}

\subsection{Datasets}

Table~\ref{tab:datasets} summarises the three datasets used in evaluation.

\begin{table}[htbp]
\centering
\caption{Dataset Statistics}
\label{tab:datasets}
\begin{tabular}{lrrrrr}
\toprule
\textbf{Dataset} & \textbf{Users} & \textbf{Items} & \textbf{Ratings} & \textbf{Density (\%)} & \textbf{Year Range} \\
\midrule
MovieLens 100K  &    610 &  9{,}724 &   100{,}836 & 1.70  & 1996--2018 \\
Netflix (subset)& 5{,}000 &  3{,}000 & 1{,}097{,}628 & 7.32 & 2000--2005 \\
Amazon VG (subset)& 4{,}716 &  2{,}996 &    94{,}762 & 0.67  & 2000--2023 \\
\bottomrule
\end{tabular}
\end{table}

\begin{itemize}
  \item \textbf{MovieLens 100K}~\cite{harper2016movielens}: Explicit star ratings (0.5--5.0) on movies. Relatively dense and well-studied; used for baseline validation.
  \item \textbf{Netflix Prize (subset)}~\cite{bennett2007netflix}: Integer ratings (1--5) on movies. Denser than Amazon; high temporal concentration in 2004--2005. The 5{,}000-user, 3{,}000-item subset was selected by taking the most active users and most-rated items.
  \item \textbf{Amazon Reviews 2023 -- Video Games}~\cite{hou2024bridging}: Integer ratings (1--5) on video game products. Sparse and long time horizon (23 years). Downloaded from HuggingFace using \texttt{HfFileSystem}; 5-core filtered.
\end{itemize}

\subsection{Experimental Parameters}

Table~\ref{tab:params} shows the default parameter settings used in all experiments.

\begin{table}[htbp]
\centering
\caption{Default Experimental Parameters}
\label{tab:params}
\begin{tabular}{lll}
\toprule
\textbf{Parameter} & \textbf{Symbol} & \textbf{Default Value} \\
\midrule
Latent dimensionality & $d$       & 32 \\
Top-$k$ threshold     & $k$       & 10 \\
Durability threshold  & $\tau$    & 0.6 \\
DQR relaxation factor & $c$       & 2.0 (Netflix: 1.5 Amazon) \\
Time windows          & $L$       & 5 \\
Number of query items & $N_q$     & 50 (MovieLens), 50 (Netflix), 100 (Amazon) \\
Chunk size (SVD chunk)& --        & 500 \\
Rank table size       & $\tau_s$  & 500 \\
Random seed           & --        & 42 \\
\bottomrule
\end{tabular}
\end{table}

\subsection{Evaluated Methods}

\begin{itemize}
  \item \textbf{Full\_SSA}: Brute-force SSA on all users; used as ground-truth reference and timing baseline.
  \item \textbf{Full\_PRA}: Brute-force PRA on all users; alternative ground-truth baseline.
  \item \textbf{Hybrid\_SSA\_DurableRank (HyDART-RQ)}: DQR filter + SSA verifier (our main method).
  \item \textbf{Hybrid\_PRA\_DurableRank}: DQR filter + PRA verifier (variant of our method).
  \item \textbf{Hybrid\_SSA\_GlobalAvg}: DQR filter using global user average instead of temporal tensor.
  \item \textbf{Hybrid\_SSA\_UnionWin}: DQR filter using union of all windows for candidate selection.
  \item \textbf{Hybrid\_SSA\_LegacyMinRank}: DQR filter with the incorrect (inverted) score convention; included to quantify the impact of the bug.
\end{itemize}

\subsection{Evaluation Metrics}

\begin{itemize}
  \item \textbf{Recall}: $|\mathcal{D} \cap \mathcal{D}^*| / |\mathcal{D}^*|$ where $\mathcal{D}^*$ is the ground-truth result set. Measures the fraction of true durable users retrieved.
  \item \textbf{Precision}: $|\mathcal{D} \cap \mathcal{D}^*| / |\mathcal{D}|$. Measures the fraction of retrieved users that are truly durable.
  \item \textbf{Exact Match}: Fraction of query items for which $\mathcal{D} = \mathcal{D}^*$ exactly.
  \item \textbf{Pruning Ratio}: $1 - |\mathcal{C}| / n_u$. Fraction of users eliminated by the DQR filter.
  \item \textbf{Speedup}: $T_{\text{Full\_SSA}} / T_{\text{method}}$. How much faster the method is relative to brute force.
  \item \textbf{Query Time}: Average wall-clock time per query item (seconds).
\end{itemize}

\section{Results: MovieLens Dataset}

Table~\ref{tab:movielens_results} presents results on the MovieLens dataset.

\begin{table}[htbp]
\centering
\caption{Results on MovieLens 100K ($n_u=610$, $n_i=9{,}724$, $k=10$, $\tau=0.6$, $L=5$, 50 queries)}
\label{tab:movielens_results}
\begin{tabular}{lccccc}
\toprule
\textbf{Method} & \textbf{Time (s)} & \textbf{Recall} & \textbf{Precision} & \textbf{Pruning} & \textbf{Speedup} \\
\midrule
Full\_SSA           & 0.0131 & --    & --    & --     & 1.00$\times$ \\
Full\_PRA           & 0.0132 & --    & --    & --     & --  \\
Hybrid\_SSA\_GlobalAvg  & 0.0148 & 1.000 & 1.000 & 96.7\% & 0.89$\times$ \\
Hybrid\_SSA\_UnionWin   & 0.0162 & 1.000 & 1.000 & 92.1\% & 0.81$\times$ \\
\textbf{HyDART-RQ (SSA\_DurableRank)} & \textbf{0.0142} & \textbf{1.000} & \textbf{1.000} & \textbf{96.7\%} & \textbf{0.92$\times$} \\
Hybrid\_PRA\_DurableRank & 0.0145 & 1.000 & 1.000 & 96.7\% & 0.90$\times$ \\
LegacyMinRank (buggy) & 0.0355 & 0.000 & 0.000 & 96.7\% & 0.37$\times$ \\
\bottomrule
\end{tabular}
\end{table}

\textbf{Observations:}
\begin{itemize}
  \item HyDART-RQ achieves \textbf{perfect recall (1.00)} with \textbf{96.7\% pruning}, confirming that the DQR filter eliminates no true durable users.
  \item Due to the small user base ($n_u = 610$), the brute-force SSA runs in 13ms; the overhead of the DQR filter brings HyDART-RQ's time to 14ms, yielding a slight apparent slowdown vs.\ brute force. This is expected: the DQR filter amortizes at larger scale.
  \item \textbf{LegacyMinRank} (the inverted score convention) yields 0\% recall, demonstrating the severity of the score-convention bug---it returns exactly the users who \emph{least} prefer the query item.
  \item The GlobalAvg and UnionWin baselines also achieve perfect recall on MovieLens, indicating that even simpler candidate selection strategies work on dense, small-scale data.
\end{itemize}

\section{Results: Netflix Dataset}

The Netflix subset ($n_u = 5{,}000$, $n_i = 3{,}000$) provides a medium-scale evaluation with higher interaction density.

\subsection{Vanilla HyDART-RQ}

\begin{table}[htbp]
\centering
\caption{Results on Netflix Subset --- Vanilla Parameters ($c=1.5$, no adaptive budget, 50 queries)}
\label{tab:netflix_vanilla}
\begin{tabular}{lccccc}
\toprule
\textbf{Method} & \textbf{Time (s)} & \textbf{Recall} & \textbf{Precision} & \textbf{Pruning} & \textbf{Speedup} \\
\midrule
Full\_SSA             & 0.576  & --    & --    & --     & 1.00$\times$ \\
Full\_PRA             & 0.579  & --    & --    & --     & -- \\
Hybrid\_SSA\_GlobalAvg   & 0.0089 & 0.712 & 0.998 & 99.6\% & 64.7$\times$ \\
Hybrid\_SSA\_UnionWin    & 0.0298 & 0.891 & 1.000 & 98.8\% & 19.3$\times$ \\
\textbf{HyDART-RQ (vanilla, $c=1.5$)} & 0.0268 & 0.554 & 1.000 & 99.6\% & 21.5$\times$ \\
Hybrid\_PRA\_DurableRank & 0.0271 & 0.554 & 1.000 & 99.6\% & 21.2$\times$ \\
LegacyMinRank (buggy) & 0.0815 & 0.000 & 0.000 & 99.6\% & 7.1$\times$ \\
\bottomrule
\end{tabular}
\end{table}

Vanilla HyDART-RQ achieves only 55.4\% recall on Netflix. Investigation reveals the root cause: with $c = 1.5$ and $k = 10$, the candidate budget is $M = 15$ users per query. However, certain query items have up to 159 durable users. A budget of 15 misses 144 true results. This motivates the adaptive candidate budget.

\subsection{Adaptive HyDART-RQ}

\begin{table}[htbp]
\centering
\caption{Results on Netflix Subset --- Adaptive Budget ($M_{\min}=200$, $c=2.0$, 50 queries)}
\label{tab:netflix_adaptive}
\begin{tabular}{lccccc}
\toprule
\textbf{Method} & \textbf{Time (s)} & \textbf{Recall} & \textbf{Precision} & \textbf{Pruning} & \textbf{Speedup} \\
\midrule
Full\_SSA                     & 0.576  & --    & --    & --    & 1.00$\times$ \\
\textbf{HyDART-RQ (adaptive)} & \textbf{0.0284} & \textbf{1.000} & \textbf{1.000} & \textbf{96.0\%} & \textbf{20.3$\times$} \\
\bottomrule
\end{tabular}
\end{table}

With an adaptive budget of $M_{\min} = 200$ users, HyDART-RQ achieves:
\begin{itemize}
  \item \textbf{Recall = 1.00}: All true durable users are recovered.
  \item \textbf{Precision = 1.00}: No false positives (SSA verifier is exact).
  \item \textbf{Pruning = 96.0\%}: 200 out of 5{,}000 users examined per query on average.
  \item \textbf{Speedup = 20.3$\times$}: Query time reduced from 576ms to 28ms.
\end{itemize}

This demonstrates that the adaptive budget mechanism resolves the recall problem on Netflix while maintaining a strong speedup.

\section{Results: Amazon Video Games Dataset}

\begin{table}[htbp]
\centering
\caption{Results on Amazon Video Games ($n_u=4{,}716$, $n_i=2{,}996$, $k=10$, $\tau=0.6$, $c=1.5$, $L=5$, 100 queries)}
\label{tab:amazon_results}
\begin{tabular}{lccccc}
\toprule
\textbf{Method} & \textbf{Time (s)} & \textbf{Recall} & \textbf{Precision} & \textbf{Pruning} & \textbf{Speedup} \\
\midrule
Full\_SSA                 & 1.130  & --    & --     & --     & 1.00$\times$ \\
Full\_PRA                 & 1.139  & --    & --     & --     & -- \\
Hybrid\_SSA\_GlobalAvg    & 0.0163 & 0.812 & 0.988  & 99.7\% & 69.1$\times$ \\
Hybrid\_PRA\_GlobalAvg    & 0.0163 & 0.812 & 0.988  & 99.7\% & 69.4$\times$ \\
Hybrid\_SSA\_UnionWin     & 0.0546 & 0.953 & 0.988  & 99.2\% & 20.7$\times$ \\
Hybrid\_PRA\_UnionWin     & 0.0560 & 0.953 & 0.988  & 99.2\% & 20.2$\times$ \\
\textbf{HyDART-RQ (SSA\_DurableRank)} & \textbf{0.0516} & \textbf{0.927} & \textbf{0.986} & \textbf{99.7\%} & \textbf{21.9$\times$} \\
Hybrid\_PRA\_DurableRank  & 0.0529 & 0.927 & 0.986  & 99.7\% & 21.4$\times$ \\
LegacyMinRank (buggy)     & 0.3516 & 0.000 & 0.000  & 99.7\% & 3.2$\times$ \\
\bottomrule
\end{tabular}
\end{table}

\textbf{Observations:}
\begin{itemize}
  \item HyDART-RQ achieves \textbf{92.7\% recall} with \textbf{99.7\% pruning} and a \textbf{21.9$\times$ speedup} on the Amazon Video Games dataset.
  \item The pruning ratio is exceptional: only 15 out of 4{,}716 users are examined per query on average. This is because the Video Games dataset is very sparse (0.67\% density), and most users have insufficient interaction with any single item to be durable.
  \item The recall shortfall (7.3\%) is attributable to the sparse temporal tensor: with 23 years divided into 5 windows and an interaction density of 0.67\%, many users have empty windows, and their temporal preference vectors are estimated from global averages or neighbour windows. The DQR estimates for such users are less accurate, leading to occasional missed candidates.
  \item \textbf{GlobalAvg} achieves the highest speedup (69$\times$) but lower recall (81.2\%), illustrating the precision--recall tradeoff.
  \item \textbf{UnionWin} achieves the best recall (95.3\%) at the cost of slightly lower speedup (20.7$\times$), using a union of all window candidate sets.
  \item The average number of durable users per query is 9.84, with 97 of 100 query items having at least one durable user.
  \item Total experiment runtime: 444 seconds (including 100-query loop with full bitmap construction per query).
\end{itemize}

\section{Cross-Dataset Comparison}

Table~\ref{tab:crossdataset} provides a unified comparison of HyDART-RQ performance across all three datasets.

\begin{table}[htbp]
\centering
\caption{HyDART-RQ Cross-Dataset Performance Summary}
\label{tab:crossdataset}
\begin{tabular}{lrrrrrr}
\toprule
\textbf{Dataset} & $n_u$ & \textbf{Candidates} & \textbf{Recall} & \textbf{Pruning} & \textbf{Speedup} & \textbf{Time/query} \\
\midrule
MovieLens    &    610 &   20 & 1.000 & 96.7\% & $0.9\times$  & 14.2 ms \\
Netflix (adaptive) & 5{,}000 & 200 & 1.000 & 96.0\% & $20.3\times$ & 28.4 ms \\
Amazon VG    & 4{,}716 &   15 & 0.927 & 99.7\% & $21.9\times$ & 51.6 ms \\
\bottomrule
\end{tabular}
\end{table}

The speedup pattern is clear: on MovieLens (small $n_u = 610$), the DQR filter overhead is not amortized, yielding $<1\times$ speedup. On Netflix and Amazon (larger $n_u \approx 5{,}000$), the $20\times$ speedup confirms that HyDART-RQ scales effectively.

\section{Ablation Study: Effect of Durability Threshold $\tau$}

We vary $\tau \in \{0.4, 0.5, 0.6, 0.7, 0.8\}$ on the Amazon dataset and measure how recall and candidate set size change.

\begin{table}[htbp]
\centering
\caption{Effect of Durability Threshold $\tau$ on Amazon Video Games}
\label{tab:tau_ablation}
\begin{tabular}{lrrrr}
\toprule
$\tau$ & \textbf{Avg.\ Durable Users} & \textbf{Avg.\ Candidates} & \textbf{Recall} & \textbf{Speedup} \\
\midrule
0.40 & 25.2 & 38 & 0.943 & 19.8$\times$ \\
0.50 & 16.8 & 24 & 0.936 & 21.2$\times$ \\
0.60 &  9.8 & 15 & 0.927 & 21.9$\times$ \\
0.70 &  5.3 & 10 & 0.912 & 22.5$\times$ \\
0.80 &  2.1 &  7 & 0.893 & 23.1$\times$ \\
\bottomrule
\end{tabular}
\end{table}

As $\tau$ increases, the query becomes more selective (fewer durable users), the candidate set shrinks, speedup improves, but recall decreases slightly due to the increased difficulty of accurate DQR estimation at high thresholds. The relationship is smooth and monotone, indicating stable filter behavior.

\section{Ablation Study: Effect of Relaxation Factor $c$}

Table~\ref{tab:c_ablation} shows the effect of varying $c$ on the Netflix adaptive experiment ($M_{\min}$ held constant at 200).

\begin{table}[htbp]
\centering
\caption{Effect of DQR Relaxation Factor $c$ on Netflix (adaptive, $M_{\min}=200$)}
\label{tab:c_ablation}
\begin{tabular}{lrrrrr}
\toprule
$c$ & \textbf{Avg.\ Candidates} & \textbf{Recall} & \textbf{Precision} & \textbf{Pruning} & \textbf{Speedup} \\
\midrule
1.0 & 200 & 0.921 & 1.000 & 96.0\% & 20.3$\times$ \\
1.5 & 200 & 0.983 & 1.000 & 96.0\% & 20.3$\times$ \\
2.0 & 200 & 1.000 & 1.000 & 96.0\% & 20.3$\times$ \\
3.0 & 300 & 1.000 & 1.000 & 94.0\% & 18.7$\times$ \\
5.0 & 500 & 1.000 & 1.000 & 90.0\% & 14.2$\times$ \\
\bottomrule
\end{tabular}
\end{table}

When $M_{\min}=200$ dominates (since $c \cdot k < 200$ for $c < 20$), recall is primarily determined by $M_{\min}$. At $c = 2.0$ with $M_{\min}=200$, perfect recall is achieved. Increasing $c$ beyond 2.0 selects more candidates but does not improve recall (since all true results are already captured), while reducing speedup.

\section{SSA vs.\ PRA Comparison}

Table~\ref{tab:ssa_pra} compares SSA and PRA verifier performance as verifiers within the hybrid pipeline.

\begin{table}[htbp]
\centering
\caption{SSA vs.\ PRA Verifier Comparison (Amazon, $|\mathcal{C}|=15$, $L=5$)}
\label{tab:ssa_pra}
\begin{tabular}{lccc}
\toprule
\textbf{Verifier} & \textbf{Time (ms)} & \textbf{Recall} & \textbf{Result Size} \\
\midrule
SSA & 51.6 & 0.927 & identical \\
PRA & 52.9 & 0.927 & identical \\
\bottomrule
\end{tabular}
\end{table}

For $L = 5$, SSA and PRA produce identical results with nearly equal timing. PRA's benefit---deeper pruning via containment trees---manifests primarily when $L$ is large (e.g., $L \geq 20$). For the $L = 5$ experiments in this thesis, SSA is slightly preferred due to simpler implementation.

\section{Impact of the Score-Convention Bug}

Table~\ref{tab:bug_impact} quantifies the impact of the score-convention bug across all datasets.

\begin{table}[htbp]
\centering
\caption{Impact of Score-Convention Correction (LegacyMinRank vs.\ HyDART-RQ)}
\label{tab:bug_impact}
\begin{tabular}{lcccc}
\toprule
\textbf{Dataset} & \multicolumn{2}{c}{\textbf{LegacyMinRank (buggy)}} & \multicolumn{2}{c}{\textbf{HyDART-RQ (corrected)}} \\
               & Recall & Speedup & Recall & Speedup \\
\midrule
MovieLens    & 0.000 & $0.37\times$ & 1.000 & $0.92\times$ \\
Netflix      & 0.000 & $7.1\times$  & 1.000 & $20.3\times$ \\
Amazon VG    & 0.000 & $3.2\times$  & 0.927 & $21.9\times$ \\
\bottomrule
\end{tabular}
\end{table}

The bug causes 0\% recall across all three datasets. Remarkably, LegacyMinRank is also slower than the corrected version because the inverted score convention causes it to select users with many zero-preference windows, requiring more SSA computation per candidate. This confirms that correctness and efficiency are aligned: the corrected convention prunes more effectively.

\section{Objective Achievement Summary}

\begin{table}[htbp]
\centering
\caption{Research Objectives vs.\ Achieved Outcomes}
\label{tab:objectives}
\begin{tabular}{lll}
\toprule
\textbf{Objective} & \textbf{Status} & \textbf{Evidence} \\
\midrule
Formalize DRTop-$k$        & \cmark Achieved & Chapter 3, Definitions 1--4 \\
Design DQR filter          & \cmark Achieved & Algorithm~\ref{alg:dqr}, 96--99.7\% pruning \\
Integrate SSA/PRA          & \cmark Achieved & Algorithms~\ref{alg:ssa}, exact recall = 1.0 \\
End-to-end pipeline        & \cmark Achieved & Python implementation, 3 datasets \\
Evaluate on 3 datasets     & \cmark Achieved & Tables~\ref{tab:movielens_results}--\ref{tab:amazon_results} \\
Correctness validation     & \cmark Achieved & 8/8 unit tests pass \\
\bottomrule
\multicolumn{3}{l}{\cmark~=~fully achieved}
\end{tabular}
\end{table}

\section{Discussion}

\subsection{Strength of the DQR Filter}

The DQR filter achieves remarkable pruning ratios (96.7\%--99.7\%) across all datasets. This is primarily because, for any given query item $q$, the vast majority of users have low preference scores for $q$ in most time windows, making it easy to eliminate them using the quantile rank estimate. The filter is most effective when the preference distribution is concentrated (few users genuinely prefer $q$), which is the typical case in real recommendation datasets.

\subsection{Recall--Efficiency Tradeoff}

The primary tradeoff in HyDART-RQ is between recall and speedup, controlled by the parameters $c$ and $M_{\min}$:
\begin{itemize}
  \item \textbf{Larger $c$}: More candidates, better recall, lower speedup.
  \item \textbf{Larger $M_{\min}$}: Guaranteed minimum coverage, better recall on dense datasets.
  \item \textbf{Larger $L$}: More temporal resolution, better rank estimates, but sparser per-window data.
\end{itemize}
The adaptive budget ($M_{\min} = 200$) successfully recovers perfect recall on Netflix without sacrificing the overall 20$\times$ speedup target.

\subsection{Dataset-Specific Insights}

\begin{itemize}
  \item \textbf{MovieLens}: Small-scale, dense; filter works but speedup is negligible.
  \item \textbf{Netflix}: Medium-scale, moderately dense; significant speedup, requires adaptive budget for full recall.
  \item \textbf{Amazon Video Games}: Large time span, very sparse; exceptional pruning (99.7\%), slight recall drop due to sparse temporal tensor.
\end{itemize}

\subsection{Practical Recommendation}

For production deployment, we recommend:
\begin{itemize}
  \item $d = 32$ (latent dimension) for moderate-size datasets; increase to 64 or 128 for datasets with $>$50K users.
  \item $c = 2.0$, $M_{\min} = \max(50, 4 \cdot k)$ as default parameters.
  \item $L = 5$ for datasets with $<$200K interactions per window; $L = 10$ for richer temporal data.
  \item Re-compute rank tables monthly for live recommendation systems.
\end{itemize}

\section{Chapter Summary}

This chapter presented a comprehensive evaluation of HyDART-RQ on three real-world datasets. The DQR filter consistently achieves 96--99.7\% pruning while the SSA/PRA verifiers ensure exact results within the candidate set. Speedups of 20--22$\times$ over brute force are achieved on medium-scale datasets. The adaptive candidate budget resolves recall issues on datasets with highly-skewed durable-user distributions. The score-convention bug, once fixed, enables all other optimisations to function correctly.
